% !TEX root = ../auv-thesis.tex
% Authoritative thesis source; edit this TeX file directly.
% Legacy Markdown material is reference-only and must not overwrite this file.

\chapter{决策与控制系统}

\section{引言}\label{sec:ch04-4-1}

第 2 章给出了系统的硬件骨架，第 3 章解决了``AUV 在哪里、对自己的位置有多大把握''的问题。本章接着回答下一个问题：在已有状态估计的基础上，如何把``去哪里巡检''与``怎样安全到达''这两件事拆解成可工程化的决策与控制结构？在展开之前，必须先澄清两个被工程实践反复证明的基本认识。其一，AUV 在水平方向上是欠驱动的，它没有横向执行器，因而控制器无法直接命令横向位移，只能通过航向调节间接实现横向靠拢；其二，水下巡检任务实际上由任务调度、轨迹生成、参考生成、低层执行与安全仲裁五类语义完全不同的负载构成，它们的时间尺度、确定性程度和失败代价各不相同，把它们塞进同一个控制器既不现实也不安全。本章的全部设计都建立在这两个认识之上。

\subsection{欠驱动 AUV 运动约束分析}\label{sec:ch04-4-1-1}

本文考虑的 AUV 采用鱼雷型布局：尾部一个主推进器配合十字形舵面，水平舵与垂直舵分别调节俯仰和航向。这种构型在水平面内能够通过舵角控制航向和前进速度，但在横移方向上没有直接执行器，属于典型的欠驱动系统；垂直面内则通过水平舵控制俯仰角、间接改变深度。这一约束带来一个根本性的推论：控制器不能指望在横移方向上产生直接作用力，凡是``沿电缆中心线保持横向偏移为零''的目标，最终都必须被翻译成``调整航向使 AUV 朝中心线靠拢''。这个翻译过程，正是后文视线（LOS）制导与制导层模型预测控制存在的根本动机\cite{breivik2005guidance,fossen2015los}。

在仿真环境中，AUV 接受一个多维速度指令，经内层映射为垂直舵、水平舵与推进器转速；执行器本身具有一阶惯性，舵角与推力的实际响应滞后于指令，时间常数在百毫秒量级。这一滞后看似细节，却直接约束了预测控制的时域下限：如果预测时域过短，模型还来不及``感受到''舵效就要给出最优解，预测的意义便被掏空。据此，本文把模型预测控制的预测时域取为 2 秒，以充分覆盖一阶惯性的过渡过程。

控制对象的状态向量精简保留六个变量：

\begin{equation}
\mathbf{X} = [x,\ y,\ z,\ \psi,\ u,\ w]^{\mathsf{T}}
\label{eq:ch04-state-vector}
\end{equation}

其中 \((x,y)\) 为水平面位置，\(z\) 为深度（正向下），\(\psi\) 为航向角，\(u\) 为前向速度（surge），\(w\) 为垂向速度（heave）。横向速度（sway）不作为受控量，而是通过航向偏转间接影响水平轨迹，这正是欠驱动约束在状态向量层面的直接体现。俯仰角不作为独立状态量进入预测模型，而是按目标深度误差近似生成；这一取舍牺牲了俯仰动力学的精确性，却显著降低了优化问题的维度，使求解时间维持在控制周期可接受的范围内。

\subsection{决策与控制系统的分层解耦逻辑}\label{sec:ch04-4-1-2}

理解了欠驱动约束之后，下一个问题是如何在系统层面组织决策与控制的不同职责。海缆巡检任务需要在同一个系统里处理``去哪里巡检''\,``如何沿电缆走''\,``感知不可靠时如何降级''和``出现安全风险时如何退出''四类问题，而这四类问题的时间尺度、确定性和失败代价完全不同，硬塞进单一控制器只会让系统不可维护、故障不可定位。本文据此采用分层解耦设计，把系统自上而下划分为五层：任务调度层决定当前应处于搜索、锁定、跟踪、上浮、返航还是急停哪个阶段；轨迹生成层根据任务阶段产生参考路径与参考速度；参考生成层把参考路径转化为航向、深度、速度指令；低层控制层把参考指令落实为舵角与推力；安全仲裁层则始终监控感知、通信与执行器状态，必要时限制或回退输出。

\begin{longtable}[]{@{}
  >{\raggedright\arraybackslash}p{(\linewidth - 4\tabcolsep) * \real{0.3333}}
  >{\raggedright\arraybackslash}p{(\linewidth - 4\tabcolsep) * \real{0.3333}}
  >{\raggedright\arraybackslash}p{(\linewidth - 4\tabcolsep) * \real{0.3333}}@{}}
\caption{决策与控制系统分层职责}\label{tab:ch04-control-layers}\tabularnewline
\toprule\noalign{}
层级 & 职责 & 实现思路 \\
\midrule\noalign{}
\endhead
\bottomrule\noalign{}
\endlastfoot
任务调度层 & 决定当前处于搜索/锁定/跟踪/上浮/返航/急停哪个阶段 & 行为树（Behavior Tree） \\
轨迹生成层 & 根据任务阶段产生参考路径与参考速度 & 之字形扫描 + 电缆跟踪轨迹 \\
参考生成层 & 把参考路径转化为航向/深度/速度指令 & 视线制导 / 制导层模型预测控制 \\
低层控制层 & 把参考指令转化为舵角与推力 & 内层深度-航向自动驾驶（PID） \\
安全仲裁层 & 监控感知、通信与执行器状态，必要时限制或回退输出 & 指令守卫 + 失效保护节点 \\
\end{longtable}

这种分层在工程上带来两个直接好处：一是允许各层独立验证，任务调度可以在不连接控制器的情况下用模拟动作验证调度逻辑，控制器也可以脱离行为树单独做控制扫描；二是允许各层独立替换，用 PID 替换模型预测控制不影响行为树，用行为树替换状态机也不影响控制层。任务运行阶段及其准入、回退关系如图~\ref{fig:ch04-mission-state-machine}所示。该状态机描述的是任务生命周期的阶段管理；每个阶段内部的行动选择仍由后文的行为树承担，二者在职责上互补而非替代。

\begin{figure}[htbp]
\centering
\includegraphics[width=0.95\linewidth,height=\textheight,keepaspectratio,alt={IDLE、PREFLIGHT、SHADOW、SINGLE LOOP和FULL AUTONOMY的分级释放、任务完成与安全回退，任一阶段保留回退通道}]{architecture/auv_v2_mission_state_machine.pdf}
\caption{巡检任务生命周期状态机：IDLE \(\rightarrow\) PREFLIGHT \(\rightarrow\) SHADOW \(\rightarrow\) SINGLE LOOP \(\rightarrow\) FULL AUTONOMY 的分级释放、任务完成与安全回退；任一阶段保留回退通道}\label{fig:ch04-mission-state-machine}
\end{figure}
\FloatBarrier

\section{基于行为树的任务规划与调度}\label{sec:ch04-4-2}

确定分层架构之后，第一个需要具体落实的是任务调度层。工程实践中可选的实现方式主要有两类，传统的有限状态机与现代的行为树。有限状态机直接、易写、易理解，但对``应急优先级抢占''和``任务可组合''的支持极为薄弱：每新增一种故障类型，都要在所有相关状态之间补画新的转移边，状态数很快膨胀到平方量级。行为树则把``组合''作为第一等公民，应急行为可以以高优先级子树的形式插入到任意位置而不破坏主任务逻辑\cite{colledanchise2018behavior,colledanchise2017how,sprague2018improving}。本文据此选择行为树作为任务调度层的实现，其``安全优先、任务后置''的结构如图~\ref{fig:ch04-behavior-tree-core}所示。

\begin{figure}[htbp]
\centering
\includegraphics[width=0.95\linewidth,height=\textheight,keepaspectratio,alt={Selector以安全监督优先抢占任务序列，Condition与Action构成执行关系并保留回落路径}]{architecture/auv_v2_behavior_tree_illustration.pdf}
\caption{行为树决策核心：Selector 以安全监督优先抢占任务序列，Condition 与 Action 构成条件检查和动作执行关系，并保留失败回落路径}\label{fig:ch04-behavior-tree-core}
\end{figure}
\FloatBarrier

\subsection{行为树节点设计}\label{sec:ch04-4-2-1}

行为树由控制节点、条件节点和行动节点组成：条件节点检查系统状态并返回成功、失败或运行中，行动节点执行具体操作，控制节点则按既定语义组合子节点。本文采用``选择 + 顺序''组合来实现``应急优先 + 主线任务''的优先级抢占模型，根节点为选择节点，第一子节点是应急顺序子树（电池、漏水、失联检查加紧急行动），第二子节点是主任务顺序子树（搜索、锁定、跟踪）。选择节点按从左到右的顺序检查子节点，只要应急子树中任一条件被触发，主任务就立即被抢占。这种``优先级即顺序''的设计让代码层面不需要任何显式的中断或信号机制，所有抢占语义都由节点的排列顺序直接表达。任务调度所依赖的关键条件与行动节点如下：

\begin{longtable}[]{@{}
  >{\raggedright\arraybackslash}p{(\linewidth - 4\tabcolsep) * \real{0.3333}}
  >{\raggedright\arraybackslash}p{(\linewidth - 4\tabcolsep) * \real{0.3333}}
  >{\raggedright\arraybackslash}p{(\linewidth - 4\tabcolsep) * \real{0.3333}}@{}}
\caption{行为树关键条件与行动节点}\label{tab:ch04-bt-nodes}\tabularnewline
\toprule\noalign{}
\begin{minipage}[b]{\linewidth}\raggedright
节点类型
\end{minipage} & \begin{minipage}[b]{\linewidth}\raggedright
功能
\end{minipage} & \begin{minipage}[b]{\linewidth}\raggedright
职责
\end{minipage} \\
\midrule\noalign{}
\endfirsthead
\toprule\noalign{}
\begin{minipage}[b]{\linewidth}\raggedright
节点类型
\end{minipage} & \begin{minipage}[b]{\linewidth}\raggedright
功能
\end{minipage} & \begin{minipage}[b]{\linewidth}\raggedright
职责
\end{minipage} \\
\midrule\noalign{}
\endhead
\bottomrule\noalign{}
\endlastfoot
条件 & 电池健康检查 & 检查电池电压是否高于安全阈值 \\
条件 & 电缆检出确认 & 检查声磁观测是否确认电缆存在 \\
条件 & 通信状态检查 & 检查与上位机/执行层的通信状态 \\
条件 & 安全裕度检查 & 检查离底高度是否低于安全阈值 \\
行动 & 之字形搜索 & 执行之字形扫描搜索 \\
行动 & 电缆跟踪 & 沿电缆中心线跟踪 \\
行动 & 紧急上浮 & 紧急上浮到安全深度 \\
行动 & 返航 & 返回起始位置 \\
\end{longtable}

\subsection{巡检任务逻辑流}\label{sec:ch04-4-2-2}

上一节给出了节点设计，本节回答这些节点如何被组织成具体的巡检流程，以及为什么要把巡检切成搜索、锁定、跟踪与异常处理四个阶段，而非直接做端到端跟踪。答案来自对真实巡检场景的两点观察：其一，AUV 入水后的初始位置不一定恰好在电缆正上方，海图先验可能存在数十米量级的偏差，必须先用扫描模式扩大可探测范围；其二，从``看见电缆信号''到``获得稳定的相对位置估计''之间存在一个必经的过渡过程，观测刚出现时往往噪声大、置信度低，直接进入跟踪会导致控制器频繁修正错误目标。这两点决定了``探测能力的逐步建立''必须被显式建模为不同阶段，而不能压缩成单一控制环。据此，巡检任务被组织为四个互斥但可抢占的阶段。

搜索阶段在未知区域执行之字形扫描以扩大探测范围，等待声磁观测在连续多帧内同时报告高置信度后退出；锁定阶段把粗略的电缆位置估计逐步收敛到稳定的相对位置估计上，期间通常压低航速以让滤波器在低动态条件下完成观测融合，待相对位置估计的协方差降到跟踪阶段允许的阈值以下即退出；跟踪阶段沿电缆路由持续更新航向、速度和离底高度，是巡检任务的主线工作模式，绝大部分有效观测和有效作业里程都在此阶段产生，其退出条件包括任务完成、感知质量退化触发降级、以及安全约束被违反触发应急三类；异常处理则不是独立的工作模式，而是其他三阶段在异常条件下的临时退出路径，感知不可靠时降速或退回搜索，安全约束被违反时触发紧急上浮，处理后可根据恢复情况返回相应正常阶段或进入应急子树。

四阶段之间的转换由关键观测信号驱动而非时间表驱动：从搜索切到锁定依赖声磁检测的高置信度持续帧数，从锁定切到跟踪依赖协方差的下降，从跟踪退到异常处理依赖感知健康状态的变化。把转换条件设计成``观测驱动''而非``时间驱动''的好处是，任务流程对场景的快慢变化具有自适应能力，不会因预设时间表与实际节奏不符而陷入僵硬执行。行为树相对有限状态机的核心优势在此进一步显现：四阶段之间的转换关系不需要画成平方量级的转移边，而只需在顺序节点的子节点之间按序排列条件检查，转换语义自动由节点返回值的传播实现。

\subsection{应急自救行为的优先级机制}\label{sec:ch04-4-2-3}

正常任务流程之上，行为树还必须回答最后一个问题：当多种故障同时出现时，AUV 应当按什么顺序响应？这一顺序不能由编程偏好决定，而必须从``故障的不可逆性 + 响应的时间紧迫性''两个维度推出。应急子树对低电压、漏水、失联、传感器失效和控制器异常保持比所有正常任务更高的优先级，且五类故障内部也存在严格排序：漏水 \textgreater{} 失联 \textgreater{} 低电压 \textgreater{} 控制器异常 \textgreater{} 传感器失效。这一排序的依据是，漏水一旦发生即对舱内电气与电池构成不可逆的损坏威胁，必须立即退出正常任务；失联意味着控制权链路失效，系统应停止继续释放自主指令并交由本地安全层兜底；低电压虽紧迫但仍允许``返航式''的渐进退出；控制器异常可由内层 PID 基线或零指令回退接管而不影响安全裕度；传感器失效中最常见的短时 DVL 丢包则只需降速等待恢复。换言之，优先级排序的本质是``哪类故障留给系统的反应时间最短''，而不是``哪类故障最常见''。

\begin{longtable}[]{@{}
  >{\raggedright\arraybackslash}p{(\linewidth - 4\tabcolsep) * \real{0.3333}}
  >{\raggedright\arraybackslash}p{(\linewidth - 4\tabcolsep) * \real{0.3333}}
  >{\raggedright\arraybackslash}p{(\linewidth - 4\tabcolsep) * \real{0.3333}}@{}}
\caption{应急故障触发条件与响应行为}\label{tab:ch04-emergency-priority}\tabularnewline
\toprule\noalign{}
\begin{minipage}[b]{\linewidth}\raggedright
故障类型
\end{minipage} & \begin{minipage}[b]{\linewidth}\raggedright
触发条件
\end{minipage} & \begin{minipage}[b]{\linewidth}\raggedright
响应行为
\end{minipage} \\
\midrule\noalign{}
\endfirsthead
\toprule\noalign{}
\begin{minipage}[b]{\linewidth}\raggedright
故障类型
\end{minipage} & \begin{minipage}[b]{\linewidth}\raggedright
触发条件
\end{minipage} & \begin{minipage}[b]{\linewidth}\raggedright
响应行为
\end{minipage} \\
\midrule\noalign{}
\endhead
\bottomrule\noalign{}
\endlastfoot
漏水 & 漏水传感器触发 & 立即上浮并急停 \\
失联 & PC 命令、遥测或 Jetson 链路按各自看门狗超时 & 停止释放自主指令；ROS2 回零指令，PC104 切遥控清零 \\
低电压 & 电池电压低于阈值 & 上浮到安全深度并返航 \\
控制器异常 & 求解超时或发散 & 回退到 PID 基线 \\
传感器失效 & DVL/声呐/磁场连续异常 & 减速并切换到搜索模式 \\
\end{longtable}

这五类故障的检测路径在设计上彼此独立，体现了``安全检测冗余''的工程原则：电池电压采用``连续多帧低于阈值''的滤波以避免单帧噪声误触发；漏水信号则采用``任一通道触发即报警''的或运算且不做任何滤波，漏水检测宁可误报也不可漏报；通信失联不采用单一阈值，而是由 ROS2 仲裁层、自主守卫层和 PC104/VxWorks 执行层分别检查 PC 原始命令、上行遥测新鲜度与 Jetson 模式指令有效性；DVL、声呐、磁场异常各自按健康计数累积到阈值后触发降级，深度越限、近底或 DVL 失锁还可由 PC104 本地应急逻辑覆盖为上浮或防触底动作；控制器异常则直接由求解器的返回状态捕获。任何一条检测路径出问题都不会波及其余，这正是应急自救可靠性的工程基础。当任一安全条件被违反时，行为树中断当前任务并切换到相应的应急行为，各守卫触发后的仲裁优先级和目标安全状态见第 2 章图~\ref{fig:ch02-emergency-arbitration}。为避免把架构差异与阈值整定混为一谈，行为树和有限状态机共享同一组施密特触发、连续样本确认与紧急故障直通规则；第~\ref{sec:ch04-bt-fsm-figures} 节给出同输入下的等价对照。

\subsubsection{行为树 vs 状态机对比图组}\label{sec:ch04-bt-fsm-figures}

下列实验分别从反应时延、抖动行为、复杂度规模、扩展开销和组合故障生存率五个维度对比两种方案。漏水故障在两种架构中均保持不滤波直通；1000 次持续故障注入得到相同的 100 ms 响应延迟和零遗漏，因而反应时延不构成架构差异。对置信度切换则采用完全相同的中心阈值 0.7、迟滞带宽 0.08 和 3 帧连续确认，避免用不同规则制造架构优势。在 100 条、每条 30 s 的阈值噪声序列中，单阈值方案平均发生 149.09 次模式切换，施密特触发与去抖后降至 2.23 次，降幅约 98.5\%；500 次低电压、漏水和近底穿越组合故障中，两种架构均无遗漏或死锁，任务生存率一致。因此，实验不支持``行为树响应更快''，其优势应限定为高优先级安全逻辑的集中表达和新增故障时较低的维护开销，下面仅保留能体现该优势的复杂度与扩展开销两图；反应时延、生存率与去抖降幅（98.5\%）的等价性已在上文以数值给出，不再单独作图。

\begin{figure}[htbp]
\centering
\includegraphics[width=0.95\linewidth,height=\textheight,keepaspectratio,alt={行为树与有限状态机的结构复杂度对比（圈复杂度与节点规模）}]{experiments/decision_bt_vs_fsm/03_complexity_comparison.png}
\caption{行为树与有限状态机的结构复杂度对比（圈复杂度与节点规模）}\label{fig:ch04-bt-fsm-complexity}
\end{figure}
\FloatBarrier

\begin{figure}[htbp]
\centering
\includegraphics[width=0.95\linewidth,height=\textheight,keepaspectratio,alt={新增故障分支的扩展开销对比（有限状态机需补转移边，行为树仅需插入高优先级子树）}]{experiments/decision_bt_vs_fsm/04_expansion_cost.png}
\caption{新增故障分支的扩展开销对比（有限状态机需补转移边，行为树仅需插入高优先级子树）}\label{fig:ch04-bt-fsm-expansion}
\end{figure}
\FloatBarrier

\section{不确定感知下的轨迹生成策略}\label{sec:ch04-4-3}

任务调度层决定了``当前应该做什么''，下一个问题是``按什么样的几何路径去做''。这一问题可以进一步拆成两个互补的子问题：其一，搜索阶段需要一条能够覆盖未知区域的扫描轨迹；其二，跟踪阶段在感知质量发生变化时，需要一种能在``激进巡检''与``保守等待''之间动态切换的航速规划机制。前者关注空间几何，后者关注时间动态，两者共同决定 AUV 在不同感知条件下走出什么样的物理路径。

\subsection{对齐行业标准的之字形空间扫描轨迹设计}\label{sec:ch04-4-3-1}

之字形扫描是海底巡检常用的预扫描模式\cite{choset1998boustrophedon,yordanova2020coverage}，其关键设计量不是单独的声呐波束角，而是声、磁两类传感器均能维持有效观测的联合扫描宽度。设 AUV 离底高度为 \(h\)，声呐横向波束角为 \(\beta\)，则理想几何声呐带宽为 \(W_{\mathrm{s}}=2h\tan(\beta/2)\)；将磁传感器在给定检测判据下的有效横向宽度记为 \(W_{\mathrm{m}}\)，联合扫描采用
\begin{equation}
W_{\mathrm{eff}}=\min(W_{\mathrm{s}},W_{\mathrm{m}}),\qquad
d=\lambda W_{\mathrm{eff}},\qquad
\eta=1-\lambda ,
\label{eq:ch04-coverage-spacing}
\end{equation}
其中 \(d\) 为相邻航线间距，\(\lambda\) 为间距比，\(\eta\) 为名义重叠率。取较小宽度可以防止某一通道的较大视场掩盖另一通道的覆盖缺口。

这里需要区分两类之字形：本节讨论 SEARCH 阶段的大范围覆盖扫描；TRACK / REACQUIRE 阶段则采用局部受控横切来补充磁几何可观测性。后者利用峰值时刻定位横向中心，并以半高全宽反演垂直距离，其几何作用见图~\ref{fig:ch02-active-crossing}，本节不重复插图。

为把航迹偏差纳入间距设计，设横向跟踪误差界为 \(e=\epsilon W_{\mathrm{eff}}\)。相邻航线向不利方向偏移时，每条扫描带可保证的宽度缩减为 \(W_{\mathrm{g}}=\max(W_{\mathrm{eff}}-2e,0)\)，由此得到归一化保证覆盖率
\begin{equation}
C_{\mathrm{g}}
=\min\left(1,\frac{W_{\mathrm{g}}}{d}\right)
=\min\left(1,\frac{\max(1-2\epsilon,0)}{\lambda}\right).
\label{eq:ch04-guaranteed-coverage}
\end{equation}
该式把``两成重叠裕度''转化为可检验条件：当横向误差界为有效宽度的 \(10\%\) 时，\(\lambda=0.8\) 恰好仍满足 \(C_{\mathrm{g}}=1\)。

覆盖安全性还需与航程和时间共同评价。对沿航线方向长度为 \(L\)、横跨方向宽度为 \(B\) 的矩形区域，航线数取
\[
n=\left\lceil\frac{B-W_{\mathrm{eff}}}{d}\right\rceil+1 ,
\]
并以直线扫描段加直径为 \(d\) 的半圆换向近似总航程 \(S=nL+(n-1)\pi d/2\)，任务时间为 \(T=S/v\)。本文以 \(W_{\mathrm{eff}}\) 的参考尺度归一化区域、航程和速度，对声呐波束角、磁有效宽度、间距比、速度比及横向误差比执行了 \(4\times4\times5\times4\times4=1280\) 组全因子解析扫描。表~\ref{tab:ch04-normalized-coverage}给出参考波束、等宽声磁通道和 \(\epsilon=0.10\) 下的间距对照。

\begin{longtable}[]{@{}rrrrrr@{}}
\caption{横向误差比为 0.10 时的归一化覆盖与时间权衡}\label{tab:ch04-normalized-coverage}\tabularnewline
\toprule\noalign{}
间距比 & 重叠率 & 保证覆盖率 & 航线数 & 归一化航程 & 时间指数 \\
\midrule\noalign{}
\endfirsthead
\toprule\noalign{}
间距比 & 重叠率 & 保证覆盖率 & 航线数 & 归一化航程 & 时间指数 \\
\midrule\noalign{}
\endhead
\bottomrule\noalign{}
\endlastfoot
0.6 & 40\% & 100.0\% & 33 & 1350.16 & 1.311 \\
0.7 & 30\% & 100.0\% & 29 & 1190.79 & 1.156 \\
0.8 & 20\% & 100.0\% & 25 & 1030.16 & 1.000 \\
0.9 & 10\% & 88.9\% & 23 & 951.10 & 0.923 \\
1.0 & 0\% & 80.0\% & 20 & 829.85 & 0.806 \\
\end{longtable}

\begin{figure}[htbp]
\centering
\includegraphics[width=0.88\linewidth,height=\textheight,keepaspectratio,alt={横向误差比为百分之十时扫描间距、保证覆盖率与任务时间的归一化权衡}]{experiments/coverage/normalized_coverage_time_tradeoff.png}
\caption{扫描间距对保证覆盖率与归一化任务时间的影响}\label{fig:ch04-normalized-coverage-tradeoff}
\end{figure}
\FloatBarrier

结果表明，\(\lambda=0.8\) 是当前离散设计集中仍保持完全保证覆盖的最大间距；增至 0.9 后保证覆盖率降为 88.9\%，而减至 0.7 和 0.6 虽继续保持完全覆盖，却分别增加约 15.6\% 和 31.1\% 的任务时间。速度只按反比改变任务时间而不改变几何覆盖，波束角与磁有效宽度则通过二者的较小值共同决定航线数量。因此，0.8 倍联合有效宽度不是孤立经验常数，而是在 10\% 归一化横向误差界下的覆盖--时间折中点。这里给出的均为无量纲设计规律；实际米制间距与任务时长仍须由实测波束图、作业高度、磁检测宽度、最小转弯半径和速度上限共同标定。

\subsection{基于感知置信度的动态航速规划模型}\label{sec:ch04-4-3-2}

空间扫描的几何确定之后，下一个问题是时间维度上的速度调节。在巡检阶段，AUV 不应当总以恒定速度前进，感知置信度变化时，行驶节奏也应随之变化。一个朴素的直觉是：感知质量稳定时可以按设计速度沿电缆巡检，感知置信度下降时则应降低航速、增加安全裕度。把这一直觉量化的核心，是将感知不确定性映射为速度上限的连续调节函数，而不是仅在``正常/异常''两种离散状态之间硬切换。

为什么不用硬切换而要用连续映射？原因有二。其一，硬切换会在阈值附近产生抖动，当置信度在阈值上下小幅波动时，速度上限会反复跳变，导致 AUV 不停加减速；其二，连续映射允许``轻度退化时小幅减速、严重退化时大幅减速''，把可用的感知信息按比例利用，而不是把所有不可靠观测粗暴地一律降级。这一机制与后文的不确定性感知控制权重调整深度耦合，其运作可归纳为三点。

首先是感知置信度到速度上限的映射。系统维护一个实时置信度，由状态估计的协方差经映射模块换算得到（取值介于零与一之间）；当置信度低于设定阈值时，参考速度随置信度下降而降低，形成``感知越不可靠、行驶越保守''的行为。其次是与控制权重联动的闭环：置信度下降不仅压低速度上限，同时触发控制代价权重的动态调整，增大控制惩罚使指令更平滑，放大跟踪权重使 AUV 在已知参考附近保守跟随。速度通道决定``走多快''，权重通道决定``走多稳''，两者协同地把``感知不确定''翻译成 AUV 物理行为上的双重保守。最后是降级策略的层次化响应：置信度轻度下降时首先做速度调节，中度下降时切换回搜索模式，严重下降时触发紧急上浮。这三级响应不是相互替代，而是随置信度持续恶化按梯度逐级触发，先减速等待、再退回搜索、最后上浮，为系统提供从轻度退化到严重退化的连续应对手段。整条链路可概括为：状态估计协方差经映射得到置信度，再同时驱动控制代价权重与参考速度，其协方差到置信度的映射细节见第 3 章。

\section{不确定性感知模型预测控制设计}\label{sec:ch04-4-4}

任务调度层与轨迹生成层就位之后，本章进入最关键的一层，也就是本文最主要的算法贡献：不确定性感知的模型预测控制（Uncertainty-Aware MPC，下称 UA-MPC）。在具体构造上，其约束设置与制导层定位受到约束型 MPC 和 AUV 轨迹跟踪研究的启发\cite{mayne2000constrained,li2021trajectory}；区别在于，本文进一步把在线感知置信度引入代价权重。进入具体公式之前，需要先澄清一个工程上的定位问题：本系统中的 MPC 并不是直接驱动舵机的低层控制器，而是\textbf{制导层的参考生成器}。作出这一定位有两条理由。其一，内层已经有一个调优良好的深度-航向 PID 自动驾驶，强行用 MPC 替代它不仅得不到额外收益，反而引入不必要的求解时间；其二，MPC 的真正优势在于多变量耦合、约束处理、未来预瞄和不确定性权重，只有在制导层才能充分体现，而非在单一深度闭环里。本节按``模型、代价、不确定性权重、参考、约束、求解器''六个层次依次展开。

\subsection{离散化运动学预测模型}\label{sec:ch04-4-4-1}

按上述路线，本节首先回答``模型层''应当采用什么样的状态空间表示。MPC 是一种模型驱动的控制方法，预测模型是它的全部物理依据，模型设计是否合理直接决定后续代价函数与约束的讨论是否有意义。这里从两个问题入手：状态向量应取多少维？离散化应采用哪种数值方案？

先看状态维度的选择。全自由度刚体模型把运动分解为三个平动与三个转动，连同位形和速度共十二维。但本系统的 MPC 不需要承担这一全套，理由有三：其一，欠驱动 AUV 在水平面没有横向执行器，横向速度由航向偏转间接产生，把它作为独立状态量不会带来控制自由度，反而使求解维度翻倍；其二，横滚在巡检任务中本就被内层的小角度控制抑制，纳入预测模型只会让动力学方程复杂化；其三，俯仰角由目标深度与当前深度之差经限幅近似生成，作为派生量进入垂向通道即可，无需独立状态。综合这三点，本系统把 MPC 状态向量精简为六维：

\begin{equation}
\mathbf{X} = [x,\ y,\ z,\ \psi,\ u,\ w]^{\mathsf{T}}
\label{eq:ch04-mpc-state}
\end{equation}

其中 \((x,y)\) 为水平面位置，\(z\) 为深度（正向下），\(\psi\) 为航向角，\(u\) 为前向速度，\(w\) 为垂向速度。这一取舍在``够用即可''与``求解时间可控''之间取得了工程平衡，既能刻画路径预瞄与深度调节的核心动力学，又把单步求解时间维持在控制周期可消化的范围。相应地，控制量精简为三维：

\begin{equation}
\mathbf{U} = [\psi_{\mathrm{cmd}},\ z_{\mathrm{cmd}},\ T_{\mathrm{cmd}}]^{\mathsf{T}}
\label{eq:ch04-mpc-control}
\end{equation}

即航向参考、深度参考与推力百分比。这再次呼应了``制导层参考生成器''的定位：MPC 输出的不是舵角和转速本身，而是三个参考量，由内层执行器在每个控制周期内消化为具体的舵面动作与转速。这种``上层只给目标、下层负责落实''的解耦既保留了调优良好的内层 PID，又把多变量耦合与约束处理交给优化器完成。

把上述状态量与控制量代入连续时间动力学，得到六个运动学方程：

\begin{equation}
\begin{aligned}
\dot{x} &= u\cos\psi, &
\dot{y} &= u\sin\psi, \\
\dot{z} &= -u\sin\theta+w\cos\theta, &
\dot{\psi} &= k_{\psi}(\psi_{\mathrm{cmd}}-\psi), \\
\dot{u} &= \frac{T_{\mathrm{actual}}-d_u u|u|}{m_u}, &
\dot{w} &= \frac{-d_w w+k_z(z_{\mathrm{cmd}}-z)+f_b}{m_w}.
\end{aligned}
\label{eq:ch04-continuous-model}
\end{equation}

其中俯仰角 \(\theta\) 由目标深度误差经比例增益限幅得到，限幅范围约为正负三十度。作这一限幅有两重工程动机：一是 AUV 实际能够稳定保持的俯仰角受控制权威与水动力恢复力矩约束，超过约三十度后舵效非线性显著加剧，预测模型赖以成立的小角度近似便会失效；二是限幅使``目标深度→俯仰角→垂向速度''这一派生链条不会因深度误差过大而产生不切实际的瞬时俯仰指令。在深度正向下的约定里，目标深度更深时对应负俯仰角；推力则经非负限幅，确保不产生反向推力，本系统主推为单向桨，物理上不允许负推力。式中其余系数（航向增益、阻力、质量、深度到垂速的耦合增益与浮力项等）的取值汇总于附录 A.4。

离散化采用显式欧拉向前法：

\begin{equation}
\mathbf{X}_{k+1}
= \mathbf{X}_k+\Delta t\,f(\mathbf{X}_k,\mathbf{U}_k)
\label{eq:ch04-euler-discretization}
\end{equation}

为什么用显式欧拉而非四阶龙格-库塔一类的高阶方法？答案来自实时控制的求解时间预算。整个预测时域内有二十步、每步六维状态与三维控制，合计数百个变量参与优化，求解器在每次迭代中都要构建并因式分解雅可比矩阵；离散化方案越复杂，单次构建时间越长。四阶方法把每步动力学求值次数由一次提高到四次，对雅可比的稀疏结构和符号微分时间都有数倍的放大效应，对当前的实时性要求是不可接受的。显式欧拉在 0.1 秒步长下的局部截断误差为步长平方量级，对 2 秒预测时域内的位置预测通常只有分米级偏差，已足以支撑制导层的参考生成；真正需要精确积分的低层执行由内层 PID 以更短步长完成，无需 MPC 自身承担。据此，时间步长取 0.1 秒、预测步数取二十步、预测时域覆盖 2 秒未来状态，这一时域选择同样不是任意的，而是由``百毫秒级一阶执行器惯性需要数倍时间常数才能完整反映稳态''反推得出，与 4.1.1 节关于执行器一阶惯性的讨论一致。

\subsection{代价函数构建}\label{sec:ch04-4-4-2}

模型层确定之后，MPC 在每个控制周期求解的优化问题，其代价由跟踪代价与控制代价两部分构成：

\begin{equation}
\min J = J_{\mathrm{tracking}}+J_{\mathrm{control}}
\label{eq:ch04-cost-decomposition}
\end{equation}

跟踪代价对预测时域内每一步累加状态与参考之差的加权平方：

\begin{equation}
\begin{aligned}
J_{\mathrm{tracking}}=\sum_k \bigl[
&w_x(x_k-x_{\mathrm{ref}})^2+w_y(y_k-y_{\mathrm{ref}})^2 \\
&+w_z(z_k-z_{\mathrm{ref}})^2+w_\psi(\psi_k-\psi_{\mathrm{ref}})^2 \\
&+w_u(u_k-u_{\mathrm{ref}})^2+w_w(w_k-w_{\mathrm{ref}})^2
\bigr].
\end{aligned}
\label{eq:ch04-tracking-cost}
\end{equation}

各维权重的相对大小反映了它们对巡检任务的重要程度：近底巡检中深度跟踪权重最大，因为深度误差直接侵蚀离底安全裕度；航向次之；水平面位置与速度则相对宽松。控制代价对控制时域内每一步的指令幅值施加二次惩罚，以抑制输出的剧烈波动：

\begin{equation}
J_{\mathrm{control}}
=\sum_k s_u(c)\left(
W_{\psi_{\mathrm{cmd}}}\psi_{\mathrm{cmd},k}^2
+W_{z_{\mathrm{cmd}}}z_{\mathrm{cmd},k}^2
+W_T T_{\mathrm{cmd},k}^2
\right)
\label{eq:ch04-control-cost}
\end{equation}

其中的整体缩放因子 \(s_u(c)\) 由不确定性感知机制根据置信度动态调整，详见 4.4.3 节。跟踪权重与控制权重的具体默认取值汇总于附录 A.4。

\subsection{不确定性感知的权重动态调整算法}\label{sec:ch04-4-4-3}

UA-MPC 是本文的核心创新之一。在给出公式之前，先把它要解决的直觉问题讲清楚：当感知可靠时，控制器应当积极跟踪参考路径、果断给出指令；当感知变得不可靠时，控制器应当自动收敛激进程度，既在已知参考附近保守跟随、不贸然远离，又为必要的纠偏动作保留控制权威，并由速率与带宽约束抑制指令抖动。传统 MPC 的代价权重是固定的，无法表达``感知质量在变、控制风格也应随之变''这层意思。已有的鲁棒 MPC 与管道（tube）MPC 通过约束收紧来抵御有界扰动\cite{mayne2005robust,jimoh2024tube}，灵敏度感知 MPC 则以对参数不确定性的灵敏度为代价加权\cite{belvedere2025sensitivity}，但二者都不直接以在线感知置信度调制代价权重。UA-MPC 的做法，就是把感知置信度引入代价权重，使跟踪与控制两类权重都随置信度连续变化。它体现在两个相互配合的机制上。

第一个机制是跟踪权重的不确定性感知放大。当置信度下降时，水平位置与航向维度的跟踪权重被放大，使优化器更倾向于贴着已知参考走、减少偏离：

\begin{equation}
\rho(c)=(1-c)^\alpha,
\qquad
w_e(c)=w_{e,0}\left[1+s_e\rho(c)\right]
\label{eq:ch04-confidence-weight}
\end{equation}

式中 \(c\in[0,1]\) 为感知置信度，\(s_e\) 为低置信度跟踪权重的附加增益，\(\alpha\) 为置信度映射的幂指数。选择幂函数 \(\rho(c)=(1-c)^\alpha\) 而不是线性映射，是为了让权重在置信度较高时几乎不变、只有当置信度明显下降时才快速放大，从而把``保守跟随''的代价集中施加在真正不可靠的区段；当幂指数取一时，该式退化为线性调整。深度维度的放大系数取水平维度的一半，因为深度通道另有内层 PID 直接兜底，不需要与水平通道同等强度的保守化。

第二个机制是控制代价的平滑调整。前一节代价函数中的整体缩放因子随置信度以 S 型（sigmoid）曲线过渡：

\begin{equation}
\begin{aligned}
\sigma(x)&=\frac{1}{1+\exp(-x)},\\
s_u(c)&=\gamma_u+(1-\gamma_u)\sigma\!\left(k(c-c_0)\right).
\end{aligned}
\label{eq:ch04-control-scale}
\end{equation}

式中 \(k\) 控制过渡的陡峭程度，\(c_0\) 为置信度转折点，\(\gamma_u\) 为低置信度下控制代价的最小缩放比例。置信度高于转折点时缩放因子接近一，控制指令正常输出；低于转折点时缩放因子平滑降到下限，控制努力惩罚随之减小，使优化器在跟踪权重已放大的同时仍保留必要的纠偏能力；指令平滑性则由后述速率与带宽约束保证。之所以用 S 型曲线而非硬阈值，是因为硬切换会在阈值附近让权重突变、进而引起控制量跳变，而平滑过渡消除了这一隐患。

为便于消融验证，UA-MPC 保留一个模式开关，可在``启用完整不确定性感知机制''与``忽略置信度、权重恒定的基线''之间切换，后者作为对照基线。围绕上述两个机制的可调量，放大上限、幂指数、S 型陡峭度、置信度阈值、控制缩放下限等，本文设计了一组消融变体，用于逐一剥离各机制的贡献。这些参数的默认取值与消融定义整理于附录 A.3。

实验章的消融从可解性和横向精度两个相互独立的指标验证了上述分解。在三场景、三种子的 PVS 闭环中，固定权重与关闭 S 型控制代价降权的变体均退化为全程回退；保留该降权后，默认 UA-MPC 与线性置信度映射变体的回退率显著下降。在确定性长波 S 弯中，默认 UA-MPC 使横向 RMSE 相对关闭 S 型降权时下降。两层证据共同说明，低置信度下减小控制努力惩罚、为纠偏保留控制权威，是可解性与精度改善的主要来源；把 \(\alpha\) 从 1.5 改为 1 仅在该运动学基准中带来极小附加改善，不足以支持修改默认参数。

\subsubsection{MPC 控制效果实验图组}\label{sec:ch04-mpc-response-figures}

图~\ref{fig:ch04-mpc-summary}以 2×2 汇总方式给出制导层 MPC 在深度阶跃、航向阶跃、电缆深度跟踪与电缆航向跟踪四类响应上的形态：四个通道均无超调发散、稳态收敛，验证了离散化预测模型与硬约束设计在接口层的可用性。单通道逐图（深度/航向阶跃、电缆深度/航向跟踪）为 PVS \(n=1\) 单元验证性质，随附录 A.7/Artifact Manifest 归档，主文不再逐图展开；控制器的多场景稳健结论由后续多场景实验统一判定。

\begin{figure}[htbp]
\centering
\includegraphics[width=0.80\linewidth,height=\textheight,keepaspectratio,alt={MPC 综合性能汇总（深度、航向与电缆跟踪四场景 2×2 汇总，PVS 仿真，\(n=1\)）}]{experiments/control_mpc/05_mpc_combined_summary.png}
\caption{MPC 综合性能汇总（深度、航向与电缆跟踪四场景 2×2 汇总，PVS 仿真，\(n=1\)）}\label{fig:ch04-mpc-summary}
\end{figure}
\FloatBarrier

\subsection{参考轨迹生成}\label{sec:ch04-4-4-4}

模型、代价、不确定性权重三层确定后，下一个问题是优化器``追什么''，也就是参考从哪里来。不同任务阶段的参考来源截然不同：跟踪阶段需要从已知电缆点序列中抽取一段沿路前推的预瞄路径，搜索阶段则没有路由可依、只能由当前航向与深度自行延展。把这两类语义压进同一段逻辑会让接口臃肿，本系统据此把参考生成分为两种互斥模式，由一个开关切换。

第一种是视线动态模式，用于已经获得稳定电缆路由的跟踪阶段。它基于电缆点序列与视线前视距离动态生成参考，算法沿电缆点序列逐步前推，前视距离随时间递增。这种``边走边伸前视''的设计意味着：预测时域内各步参考点对应的不是当前视线锚点的原地复制，而是沿电缆方向逐步向前的真实路径采样，从而让优化器可以同时兼顾``当前位置如何修正''和``未来若干步即将面对的弯道''。第二种是恒定航向模式，用于搜索与锁定阶段。它不依赖电缆点先验，完全由当前航向、当前深度和参考速度推算未来若干步的状态，使 MPC 在尚未确认电缆位置时仍能工作，而不至于因``参考缺失''而求解失败。

无论哪种模式，最终输出都是一份维度与预测时域对齐的参考矩阵，每一列对应一个时间步的期望状态、与六维状态向量一一对应。把参考从``标量序列''扩展到``完整状态序列''是 MPC 的标准做法，它让代价函数可以对每个状态分量单独施加权重，而不是只对一个综合误差求平方，这一点已在 4.4.2 节的代价函数中体现。

\subsection{执行器物理饱和约束处理}\label{sec:ch04-4-4-5}

参考准备好之后，下一个问题是优化器可以``在多大范围内''自由选择控制量。MPC 在数学上可以输出任意指令值，但实际执行器存在舵角限幅、变化率限制、推力上下限等物理约束；任何超出这些约束的指令在低层都会被强制限幅，反而使 MPC 的模型预测与真实执行轨迹脱节。本节把约束体系分为``硬约束''与``速率/带宽扩展约束''两层，前者保证指令物理可执行，后者保证指令工程可平滑。

硬约束体系直接对应执行器的物理边界：

\begin{longtable}[]{@{}lll@{}}
\caption{模型预测控制硬约束体系}\label{tab:ch04-hard-constraints}\tabularnewline
\toprule\noalign{}
约束类型 & 表达式 & 物理含义 \\
\midrule\noalign{}
\endfirsthead
\toprule\noalign{}
约束类型 & 表达式 & 物理含义 \\
\midrule\noalign{}
\endhead
\bottomrule\noalign{}
\endlastfoot
航速下限 & \(u_k\ge u_{\min}\) & 确保舵效 \\
航向指令限幅 & \(\psi_{\mathrm{cmd}}\in[\psi_{\min},\psi_{\max}]\) & 航向指令合理性 \\
深度指令限幅 & \(z_{\mathrm{cmd}}\in[z_{\min},z_{\max}]\) & 深度指令合理性 \\
推力下限 & \(T_{\mathrm{cmd}}\ge T_{\min}\) & 防止推力清零失稳 \\
推力上限 & \(T_{\mathrm{cmd}}\le T_{\max}\) & 执行器饱和 \\
\end{longtable}

各项硬约束都有明确动机：航速下限来自欠驱动 AUV 的物理事实，舵面靠水流作用力产生力矩，速度太低时舵效近乎为零，控制器再多调整也无法形成有效转向力，因此必须给航速设下限；推力下限取零而非任意负值，是因为主推为单向桨、物理上不允许反推；深度上限则与本系统设计的最大下潜深度对齐。

在硬约束之外，还需要再叠加一层速率与带宽约束：

\begin{longtable}[]{@{}
  >{\raggedright\arraybackslash}p{(\linewidth - 4\tabcolsep) * \real{0.3333}}
  >{\raggedright\arraybackslash}p{(\linewidth - 4\tabcolsep) * \real{0.3333}}
  >{\raggedright\arraybackslash}p{(\linewidth - 4\tabcolsep) * \real{0.3333}}@{}}
\caption{模型预测控制速率与带宽扩展约束}\label{tab:ch04-rate-band-constraints}\tabularnewline
\toprule\noalign{}
\begin{minipage}[b]{\linewidth}\raggedright
约束类型
\end{minipage} & \begin{minipage}[b]{\linewidth}\raggedright
表达式
\end{minipage} & \begin{minipage}[b]{\linewidth}\raggedright
物理含义
\end{minipage} \\
\midrule\noalign{}
\endfirsthead
\toprule\noalign{}
\begin{minipage}[b]{\linewidth}\raggedright
约束类型
\end{minipage} & \begin{minipage}[b]{\linewidth}\raggedright
表达式
\end{minipage} & \begin{minipage}[b]{\linewidth}\raggedright
物理含义
\end{minipage} \\
\midrule\noalign{}
\endhead
\bottomrule\noalign{}
\endlastfoot
深度指令速率 & \(\lvert z_{\mathrm{cmd},k+1}-z_{\mathrm{cmd},k}\rvert\le\Delta z_{\max}\) & 深度指令变化率 \\
航向指令速率 & \(\lvert\psi_{\mathrm{cmd},k+1}-\psi_{\mathrm{cmd},k}\rvert\le\Delta\psi_{\max}\) & 航向指令变化率 \\
深度指令带宽 & \(\lvert z_{\mathrm{cmd},k}-z_{\mathrm{current}}\rvert\le z_b\) & 深度指令不能离当前态太远 \\
航向指令带宽 & \(\lvert\psi_{\mathrm{cmd},k}-\psi_{\mathrm{current}}\rvert\le\psi_b\) & 航向指令不能离当前态太远 \\
\end{longtable}

为什么硬约束之外还要再加这一层？答案来自实测中观察到的两类不利现象。其一，纯硬约束允许优化器在相邻两步给出相差数十度的航向指令，这种跳变会让内层 PID 频繁过冲，把``平滑参考''变成``阶跃参考''；其二，优化器可能在某一步给出远离当前态的指令，例如当前深度五米、下一步却给二十余米的目标深度，即便落在深度硬限幅之内是合法的，但这种``瞬时大步跳''无法被一阶执行器跟随。速率约束限制相邻步之间的指令斜率，带宽约束限制任一时刻指令偏离当前态的最大距离，前者管短期变化率、后者管中期偏离量，两者一起把优化器的解空间收敛到内层 PID 能稳定跟随的窗口内。这两类约束可以分别启用或禁用以便消融：早期基线仅保留硬约束体系，引入速率与带宽约束后控制指令更平滑、求解时间略有上升，这一权衡由控制平滑度对照给出量化证据。上述各约束的默认取值列于附录 A.4。

\subsection{求解器实现与热启动}\label{sec:ch04-4-4-6}

约束体系给定后，最后一步是把``如何求解这个非线性优化问题''的工程实现说清楚。MPC 在每个 50 ms 控制周期都要重新求解一次包含数百个变量的非凸非线性优化问题，求解器的选择与实现直接决定整个系统能否实时运行。本系统采用符号建模框架配合内点法求解器\cite{andersson2019casadi,waechter2006ipopt}：前者把代价函数与约束自动符号微分得到雅可比，后者以内点迭代逐次逼近最优解。要做到实时可用，工程实现上有三个关键考量。

其一是``构建一次、复用多次''。优化问题的整体结构，变量数量、约束类型、雅可比稀疏模式，在配置确定后就固定，不再随时间变化。据此，本系统在初始化阶段一次性完成求解器结构的编译，运行期反复求解时只更新参考、当前状态、置信度等参数值，而不重复符号微分与编译。这一处理把单次求解时间从秒级压缩到毫秒量级，使 20 Hz 控制更新变得可实现。其二是热启动。每个周期都用上一次的最优控制序列作为本次的初始猜测，并把该序列整体向前平移一步、最后一步沿用上一步的末端控制量。这样做的物理含义是，上一周期覆盖的预测区间与本周期只相差一步，把上次解平移一步作为猜测后，求解器只需在相邻解附近继续搜索。三种子的闭环对照中，热启动使平均迭代数由 18.08 降至 9.98，平均求解时间由 17.40 ms 降至 11.29 ms。其三是求解失败的兜底。当求解器返回失败状态（超时、约束不可行或迭代发散）时，控制层捕获这一异常，并根据配置决定是保持上一次输出以维持指令连续性，还是直接退化为开环设定点输出。这条兜底路径正是 4.2.3 节``控制器异常→回退到 PID 基线''应急规则的实际接入点。

回退机制只能保证输出连续，不能替代对失败原因的观测。为区分``求解困难''、``初值不连续''和``控制周期被阻塞''，本文对每个控制周期统一记录求解器返回状态与迭代数、是否实际使用热启动、相邻状态初值跳变量、平移前后控制初值的均方根变化、初始点与最终迭代点的最大约束违反量、活跃约束数量以及当前周期墙钟时间。活跃约束按距上下界不超过 \(10^{-4}\) 判定；周期墙钟时间超过控制周期时另记阻塞事件。若失败后沿用上一成功控制指令，沿用的只是物理输出，求解状态、迭代数和耗时仍取当前失败周期，不能把上一成功帧的求解时间重复写入失败帧。由此，回退率反映``有多少周期没有产生新解''，当前墙钟时间反映``失败是否阻塞实时更新''，约束残差与初值连续性则用于定位``为什么没有产生新解''，三类语义彼此分离。

在数值容差上，本系统取一个适中的收敛容差，并把最大迭代次数设为 100，以限制单次失败对 50 ms 控制周期的占用。三场景、三种子的独立扫描表明，把上限从 100 提高到 200 仅使事件加权回退率由 0.699 降至 0.685，失败求解占用由 88.41 ms 增至 192.21 ms；提高到 400 后回退率仍为 0.594，失败占用达到 372.20 ms，诊断帧数同时减少约 26\%。因此，迭代上限偏紧只构成次要因素，继续放宽会以阻塞控制更新换取有限的求解成功率，不能作为工程修复。统一预测时域、权重和约束口径后，热启动硬约束配置在三种子闭环中的回退率与周期阻塞率均为 0.0086；关闭热启动后两者均为 0.0169。加入软约束松弛并未改善可解性，反而增加优化维度和阻塞，故默认关闭。本文据此冻结迭代上限 100、热启动开启、软约束关闭的配置；终端保持由任务层在路径完成后接管，不与求解器可解性混为一项机制。此外，指令限幅在系统中是三层冗余的：优化器内部对控制量限幅、MPC 输出后再次限幅、内层执行器最后又做一次舵角与角速度限幅，任何一层失效都不会让超界指令穿透到执行器。求解器配置与回退语义列于附录 A.5。

\subsection{近底安全过滤与控制障碍函数}\label{sec:ch04-4-4-7}

地形跟随和电缆巡检都要求车辆在尽量靠近海床的同时保持最小安全净空。单纯把目标离底高度设为常数并不足够，因为海床估计、车辆深度响应和控制周期离散化都会引入滞后：当未来数秒内海床快速抬升时，当前时刻看似安全的深度目标可能在执行到位前变成低净空风险。为此，本文在地形跟随候选深度之后叠加一层制导级安全过滤器，用控制障碍函数（control barrier function, CBF）的思想约束深度参考，而不是直接改写内层 PID 或 MPC 求解器。

令 \(z\) 表示车辆深度（向下为正），\(S\) 表示海床深度，\(d_{\min}\) 表示最小安全净空，则当前安全裕度定义为
\begin{equation}
  h = S - z - d_{\min}.
\end{equation}
当 \(h\ge 0\) 时，车辆距离海床至少保留 \(d_{\min}\) 的净空；当 \(h<0\) 时，目标深度已经越过安全边界。离散控制周期内，CBF 的核心要求不是简单把 \(h\) 保持为正，而是限制它向边界衰减的速度：
\begin{equation}
  h_{k+1} \ge (1-\alpha) h_k,\quad 0<\alpha\le 1 .
\end{equation}
这一不等式给出了一个可解释的安全缓冲：\(\alpha\) 越大，允许越快接近边界；\(\alpha\) 越小，参考深度越保守。

考虑地形跟随的预瞄需求，本文不只使用当前海床 \(S_k\)，还计算预测时间窗后的海床 \(S_{k+\ell}\)，并取二者中更浅、风险更高的一项作为过滤器的保守海床：
\begin{equation}
  S_{\mathrm{safe}}=\min(S_k,S_{k+\ell}) .
\end{equation}
候选深度 \(z_{\mathrm{ref}}\) 经过过滤后不得超过由 \(S_{\mathrm{safe}}\)、\(d_{\min}\) 和 CBF 衰减项共同给出的上界。若当前净空已经足够大，衰减项只在有限裕度内生效，避免车辆因历史大净空而长期被压在过浅深度；若当前或预瞄净空接近边界，过滤器同时限制最大下潜速率，使深度参考以车辆能够跟随的速度接近海床。

低净空状态下还需要比普通 CBF 更直接的保护。本文设置紧急净空阈值，当实际离底高度低于该阈值时，过滤器不再等待候选参考，而是强制生成上浮深度参考，并将目标航速压低；在普通接近边界但尚未进入紧急区间时，则按当前与预瞄净空裕度线性缩放 \(v_{\mathrm{ref}}\)。因此，CBF 在本文中同时作用于两个制导量：深度参考保证不继续向海床逼近，速度门控降低车辆在低净空区间内横向推进带来的响应压力。

这层设计的定位是``安全过滤''而不是``性能优化''。它不试图改善电缆横向 RMSE，也不替代跟踪授权对观测质量的判断；它只回答一个更底层的问题：在近底巡检时，候选深度是否仍位于可信安全域内。作为方法有效性的最小锚点，该过滤器在低、中、高三档地形各三种子的复核矩阵中，最小离底净空稳定保持在 2.7--2.9 m 区间，低净空违规率与穿底率均为零；完整的净空、穿底、CBF 激活与速度缩放统计及其适用边界在实验章给出。据此，本章不把该结果外推为所有极端电缆场景已安全闭合或原生执行链已经验证。

\section{决策与控制系统的鲁棒性分析}\label{sec:ch04-4-5}

前面四节给出了任务调度、轨迹生成、UA-MPC 控制与约束体系四个层面的完整设计，工程上还需要回答最后一个问题，这套设计在``非理想条件''下能否守住底线？所谓非理想条件包含两类：一是慢变退化，例如感知质量在长时间任务中持续抖动、海流扰动叠加、地形起伏比预设强烈；二是快变冲击，例如电缆突然进入掩埋段、AUV 在拐点附近遭遇横向流场、声呐通道短时完全失效。对这两类条件分别给出鲁棒性证据，是验证分层架构可工程化的关键一步。本节据此把证据组织为两条线：一条自下而上展开三层证据（近底安全→复杂路径→不确定性适应），论证各控制层级各自的能力上界；另一条自外而内收紧极端工况边界，论证哪些场景已超出当前架构的能力、必须由安全回退路径接管。

\subsection{感知信号剧烈波动下的控制平滑度对比}\label{sec:ch04-4-5-1}

本节按``近底安全→复杂路径→不确定性场景''三个层次组织鲁棒性证据，逐层说明分层控制架构在哪些条件下成立、在哪些条件下需要回退。这一顺序背后的逻辑是，任何复杂控制器在使用之前都必须先满足``不撞地''，再讨论``路径精度''，最后才谈``在不确定环境下能否保持精度''。把三层按由低到高排列，可以让读者清晰看到每层证据各自支撑哪一类工程结论，避免把``近底安全已证明''误读为``所有工况下控制器都已证明''。

第一层是近底安全性，由地形跟随基准支撑。在真口径正式结果中，PID 地形跟随实现了米级以内的离底高度跟踪，四相离底高度随真地形起伏，低、中、高三档地形消融均无低于安全阈值的违规（早先约 0.18 米的异常值系分析层基准错误造成的假象，已被真口径重跑取代，样本量与适用边界见附录 A.7）。MPC 地形跟随虽然改善了固定深度基线，但并未超过调优后的 PID 内层。这一结果意味着，在``恒定离底高度''这一最基础的近底安全任务上，调优良好的内层 PID 已把贴底误差压到米级以内，外层再加 MPC 只是锦上添花，并没有打开新的能力维度。其工程归因在于内层深度 PID 直接控制水平舵的瞬态响应，相比 MPC 制导层通过``目标深度参考加内层跟随''的两级链路少了一级延迟，对地形起伏的响应更直接。换言之，当任务的本质是``快速贴地''而非``复杂预瞄''时，MPC 的优势区被显著压缩。

第二层是复杂平面路径制导，由离线极端路径支线提供证据。在公平口径重跑中，MPC 在长波、短波 S 弯与发卡折返路径上的横向误差优于或持平于纯航向 PID 与视线制导强基线，仅在直角路径上视线制导的前瞻更优（早先``未全面超过视线制导''的判断源自实验脚手架的一处偏置缺陷，修复后已被推翻，证据边界见附录 A.7）。该结论已在注入受控横流、初始偏移与艏向噪声的五种子统计下复核，长波 S 弯与发卡掉头上 MPC 相对视线制导的优势为均值加减一倍标准差区间互不重叠。这一结果传达两层信息：一方面，预瞄能力在长波 S 弯这类连续曲率路径上把横向误差压低约一个量级（视线制导约 1.15 m 降至 MPC 约 0.13 m），这与``MPC 在预测时域内能看到下一段曲率、而纯航向 PID 只能看到当前误差''的本质差异相符；另一方面，视线制导本身已经包含某种程度的几何前视，使它在路径几何相对规则（如分段直线的直角路径）的场景下仍占优。本系统的策略由此明确，把 MPC 主要用于视线制导不擅长的不规则曲率预瞄，而不是把它当作视线制导的全面替代品。

第三层是高不确定性场景下的控制风格调制。早期主消融和 H4 横向机动对照说明，低置信度下的 S 型控制代价降权会改变求解可行域和横向机动行为；分源置信重跑后的 R13-v2 正式矩阵则给出更接近完整 ROS2 闭环的结果。由此，UA-MPC 的作用可以更准确地解释为，在观测质量变化时重新分配跟踪与控制努力，并把连续跟踪许可交给分源授权判定，而不是在感知失效后独立恢复状态；其正式矩阵的运行规模、求解耗时与控制变化率量化结果在实验章给出。最后兜底仍由 4.2.3 节定义的行为树应急规则承担。

\subsection{极端工况下的安全性边界测试}\label{sec:ch04-4-5-2}

上一节的三层证据回答了``常规工况下分层架构如何工作''，但工程验收还需要回答另一类更尖锐的问题，极端工况下分层架构在哪个边界失效，该失效是否被安全回退路径接住？为什么极端工况必须单独考虑、而不能与常规消融合并？原因有二。其一，极端工况的物理风险量级与常规场景不同：常规场景下性能下降表现为横向误差扩大几十厘米，而极端工况下相同幅度的性能下降可能直接导致撞地、绕远或丢失目标等不可逆后果；其二，极端工况之间的叠加效应不可忽略：单独的急转弯或单独的横向流都可被 MPC 的预瞄与容差吸收，但二者同时出现时其耦合行为往往不在任何单一消融场景的覆盖范围内。基于这两点，极端工况必须按``物理风险类别''与``是否触发安全回退''两个维度单独建模。

当前已有极端平面路径支线量化几何极端性，六类代理场景进一步叠加感知质量、地形与横流压力。R13-v2 正式矩阵把这类压力测试收束为四个可复核的方法侧事实：全部运行完成且无求解回退；授权模式按场景质量在 HOLD、SEARCH 与 TRACK 之间切换；保守分源置信策略显著降低控制变化率；横向 RMSE 不作为该组实验的正向主结论。综合极端另有固定权重全流程补充实验，说明固定权重基础巡缆链路在该代理场景下具备可重复闭环能力。近底安全由 CBF 三档地形三种子矩阵补充验证，各档均满足最小净空裕度、低净空违规率与穿底率为零，并保留 MPC 求解耗时上界作为实时性边界。R22 全因子离线闭环进一步把声磁量测、地形埋设与横流因子接入同一制导级闭环，覆盖 \(2^4\) 组合、多个种子与固定权重/不确定性感知两种模式，其中横流被识别为横向误差主导因子，曲线几何与横流构成主要二阶交互。上述结果共同表明，分层架构在几何、感知、地形与横流压力叠加下均由安全回退路径接住，未出现撞地或穿底。各组实验的详细数据表、终端尾段与执行链口径、以及不可外推项统一收束于实验章及其边界总表，本章不重复展开。

\subsubsection{PID/内层执行器控制效果图组}\label{sec:ch04-pid-response-figures}

图~\ref{fig:ch04-pid-summary}以四剖面汇总方式给出内层深度-航向自动驾驶在深度阶跃、航向阶跃、深度正弦与航向正弦工况下的响应形态，作为 MPC 制导层指令的内层执行基准：各剖面均稳定跟随、无持续振荡，表明内层执行接口可用。单剖面逐图为 PVS \(n=1\) 单元验证性质，随附录 A.7/Artifact Manifest 归档，主文不再逐图展开。

\begin{figure}[htbp]
\centering
\includegraphics[width=0.80\linewidth,height=\textheight,keepaspectratio,alt={内层 PID 综合剖面对比（深度、航向阶跃与正弦四剖面汇总，PVS 纯仿真复现，\(n=1\)）}]{experiments/control_pid_pvs/05_pid_pvs_profile_comparison.png}
\caption{内层 PID 综合剖面对比（深度、航向阶跃与正弦四剖面汇总，PVS 纯仿真复现，\(n=1\)）}\label{fig:ch04-pid-summary}
\end{figure}
\FloatBarrier

极端平面路径支线的量化结果与代表性轨迹图统一收纳于实验章。本章不重复展示同组结果图；其方法侧结论保持不变：MPC 在长波、短波 S 弯与发卡折返路径上优于或持平基线，但直角折弯上 LOS 的几何前瞻更优，且当前各场景均仅为单次运行。

\section{本章小结}\label{sec:ch04-4-6}

本章围绕``AUV 海缆巡检该用什么样的决策与控制结构''这一问题，按``约束→调度→轨迹→控制→鲁棒性''五个层次完成了完整设计。首先回答欠驱动 AUV 的运动约束与分层解耦的工程动机，把任务调度、轨迹生成、参考生成、低层执行、安全仲裁五类语义负载明确分层；随后给出基于行为树的任务调度方案，并通过等价规则对照把其优势限定为``故障可组合、应急可插入''的组织能力，而非快于状态机的响应速度；接着回答几何路径与速度规划，把之字形扫描与感知置信度驱动的动态航速联系起来；作为本章主线，UA-MPC 按模型、代价、不确定性权重、参考、约束和求解器六层展开，重点阐述了以幂函数放大跟踪权重、以 S 型曲线平滑控制代价的不确定性感知机制，视线动态与恒定航向两种参考模式的互斥设计，速率与带宽两类扩展约束的工程必要性，以及一次构建、热启动、失败兜底三项使求解实时可用的实现要点；在近底安全层，本文进一步给出基于预瞄海床的制导级 CBF 过滤器，把最小净空、最大下潜速率、紧急上浮和低净空降速统一为深度参考后的安全约束；最后给出三层鲁棒性证据与六类极端工况边界，明确了当前各控制层级的能力上界。

从证据等级看，本章可以稳妥支撑四类结论。内层 PID 在地形跟随主任务上保持米级以内离底误差且多档复验无安全违规，CBF 过滤器在三档地形三种子矩阵中进一步给出满足最小净空裕度、低净空违规和穿底均为零的近底安全证据；MPC 的主要优势位于需要曲率预瞄的 S 弯和发卡路径，而不是对 PID 与视线制导的全面替代；UA-MPC 的可写正结果应表述为分源置信下的控制平滑、实时可解性和跟踪许可管理，其正式矩阵无求解回退；决策层迟滞去抖在不增加紧急响应延迟的条件下显著抑制了阈值颤振。横向 RMSE、综合极端 TRACK 进入门、CBF 外推边界以及各组的具体运行数与量化指标不在本章反复展开，集中由实验章边界总表给出。

把这些结论合在一起看，本章给出的不是``一个万能控制器''，而是``一套分层控制与安全回退框架''，其工程价值不在于让某一层控制器跑得最优，而在于让任务调度、参考生成、低层执行、安全仲裁四个职责彼此独立验证、独立替换、独立兜底。这一框架将在后续实验章被置于仿真、半实物与实物链路的分层验证口径下检验，并作为硬件部署阶段的算法基线。
